% ============================================================
%  SECTION 7: ROBUSTNESS
% ============================================================

\section{Robustness}
\label{sec:robustness}

\subsection{Alternative Estimators}
\label{sec:robustness_estimators}

We compare four heterogeneity-robust difference-in-differences
estimators: Callaway--Sant'Anna
\parencite{callaway2021difference}, \textcite{sun2021estimating},
TWFE (Equation~\ref{eq:twfe}), and
\textcite{dechaisemartin2020two}. Each estimator addresses the
heterogeneous treatment effects problem through a different approach:
Callaway--Sant'Anna computes group-time ATTs relative to never-treated
units; Sun--Abraham uses interaction-weighted estimation;
de~Chaisemartin--D'Haultfoeuille applies a first-difference
estimator. TWFE is included for comparability with the prior literature
despite its known sensitivity to treatment effect heterogeneity
\parencite{goodmanbacon2021difference}.

Figure~\ref{fig:four_estimators} overlays all four estimators' dynamic
paths for both outcomes. TWFE and Sun--Abraham coefficients are
plotted relative to the omitted period ($k = -1$). Callaway--Sant'Anna
and de~Chaisemartin--D'Haultfoeuille estimates are group-time ATTs
relative to the never-treated comparison group, which do not require
a reference period. The four estimators agree on sign, magnitude, and
dynamics. Pre-treatment coefficients are statistically indistinguishable
from zero under all four estimators. Post-treatment coefficients
diverge monotonically and are statistically significant at every event
time from $k = 0$ onward.

\begin{figure}[H]
  \centering
  \includegraphics[width=\textwidth]{fig_four_estimators_combined.pdf}
  \caption{Robustness Across Four Heterogeneity-Corrected DiD Estimators}
  \label{fig:four_estimators}
  \vspace{1mm}
  \begin{minipage}{\textwidth}
    \footnotesize
    \textit{Notes:} Event-study estimates from
    \textcite{callaway2021difference}, \textcite{sun2021estimating},
    TWFE, and \textcite{dechaisemartin2020two}. TWFE and Sun--Abraham
    normalized at $k = -1$ (omitted period). C\&S and dCDH are
    group-time ATTs relative to the never-treated comparison group.
    Vertical bars denote 95\% confidence intervals. Standard errors
    clustered at the state level (27 clusters).
    Data: \textcite{chetty2020mobility}.
  \end{minipage}
\end{figure}

The TWFE static estimates (Table~\ref{tab:twfe_did}) are somewhat larger
in magnitude than the C\&S ATTs: 1.54 pp vs.\ 1.27 pp for
bottom-quintile share, and $-1.85$ pp vs.\ $-1.33$ pp for parental income rank.
This difference is consistent with TWFE placing more weight on
later post-treatment periods where the dynamic effects are larger.
The unanimity across four estimators built on fundamentally different
identifying assumptions provides strong evidence that the result is not
an artifact of any particular estimator's treatment of heterogeneous
effects.

\subsection{Timing Robustness}
\label{sec:robustness_timing}

Our preferred treatment timing assigns $t^* = 1986$ (the first birth
cohort with unambiguous full HOPE exposure). We test two alternatives:
$t^* = 1985$ (the bridge-provision cohort, potentially partially
treated) and $t^* = 1987$ (a conservative definition that discards the
partially treated cohort).

Table~\ref{tab:robustness_summary} reports static DiD estimates under
each timing definition. For bottom-quintile share, the estimates range from
1.44 pp ($t^* = 1985$) to 1.64 pp ($t^* = 1987$), all significant at
$p < 0.001$. For parental income rank, the range is $-1.78$ to $-1.85$
pp. The results are insensitive to the choice of treatment timing.

\subsection{Control Group Robustness}
\label{sec:robustness_controls}

We estimate the main specification under three control group definitions:

\begin{enumerate}
\item \textbf{Baseline}: all states without broad-based merit programs
  (26 control states, $N = 6{,}289$).
\item \textbf{Conservative}: drops additional states with any
  higher-education funding changes during the analysis window.
\item \textbf{Bordering}: restricts to states adjacent to Tennessee
  (AL, AR, GA, KY, MO, MS, NC, VA).
\end{enumerate}

Table~\ref{tab:robustness_summary} reports the results. For
bottom-quintile share, estimates range from 1.44 pp (bordering) to 1.54 pp
(baseline), all significant at $p < 0.001$. For parental income rank,
the bordering estimate ($-1.34$ pp) is somewhat attenuated relative to
the baseline ($-1.85$ pp) but remains significant.

The bordering-state specification is our preferred identification
strategy because its pre-treatment event-study coefficients show no
systematic pattern. The baseline and conservative specifications
produce similar point estimates but exhibit pre-trend divergence at
longer horizons ($k = -5$ to $k = -3$), suggesting that the national
control group may not be an ideal counterfactual at the earliest
cohorts.

\subsection{Rambachan and Roth Sensitivity}
\label{sec:robustness_rr}

We implement the honest confidence intervals of
\textcite{rambachan2023more}, which allow for bounded violations of the
parallel trends assumption. We evaluate sensitivity under two classes of restrictions.

\textbf{Relative magnitude ($\Delta_\text{RM}$).} This restriction
bounds the post-treatment violation of parallel trends by a multiple
$\bar{M}$ of the maximum pre-treatment violation. For bottom-quintile share,
the confidence interval excludes zero at $\bar{M} \leq 2$, meaning
results survive if the post-treatment trend violation is at most twice
the largest pre-treatment violation. For parental income rank, results are
robust at similar $\bar{M}$ values.

\textbf{Smoothness ($\Delta_\text{SD}$).} This restriction bounds the
change in the violation of parallel trends between consecutive periods.
Under $\Delta_\text{SD}$, the confidence intervals for both
bottom-quintile share and parental income rank include zero at moderate bounds.
This indicates fragility: if differential trends can accelerate smoothly,
the estimated effects cannot be distinguished from a continuation of
pre-existing dynamics.

The tension between $\Delta_\text{RM}$ and $\Delta_\text{SD}$ is
informative. The $\Delta_\text{RM}$ result suggests that the
post-treatment effects are large relative to pre-treatment noise,
supporting a causal interpretation. The $\Delta_\text{SD}$ result
suggests that the pre-treatment coefficients, while small, exhibit
enough curvature that a smooth extrapolation could explain the
post-treatment path. We view the $\Delta_\text{RM}$ restriction as more appropriate for this
setting. The pre-treatment coefficients are noisy but trendless---they
oscillate around zero with no systematic pattern under any of the four
estimators (Figure~\ref{fig:four_estimators}). The $\Delta_\text{RM}$
restriction bounds post-treatment violations by the magnitude of this
noise, which is the natural benchmark. The $\Delta_\text{SD}$
restriction instead bounds the \emph{rate of change} of violations,
which penalizes pre-treatment curvature even when levels are small. We
interpret the $\Delta_\text{SD}$ fragility as a genuine limitation of
the research design, but note that it reflects the sensitivity of a
particular robustness criterion rather than evidence against the
effect.

Figure~\ref{fig:rr_par_q1} displays both the $\Delta_\text{SD}$ and
$\Delta_\text{RM}$ sensitivity results for bottom-quintile share.
Figure~\ref{fig:rr_par_rank} reports the corresponding results for
parental income rank.

\begin{figure}[H]
  \centering
  \includegraphics[width=0.9\textwidth]{fig_rr_combined_par_q1.pdf}
  \caption{Rambachan--Roth Sensitivity Analysis: Bottom-Quintile Share}
  \label{fig:rr_par_q1}
  \vspace{1mm}
  \begin{minipage}{0.9\textwidth}
    \footnotesize
    \textit{Notes:} Honest confidence intervals under the smoothness
    ($\Delta_\text{SD}$) and relative magnitude ($\Delta_\text{RM}$)
    restrictions of \textcite{rambachan2023more}. The $\Delta_\text{RM}$
    restriction supports significance at $\bar{M} \leq 2$. The
    $\Delta_\text{SD}$ restriction includes zero at moderate bounds. Data:
    \textcite{chetty2020mobility}.
  \end{minipage}
\end{figure}

\begin{figure}[H]
  \centering
  \includegraphics[width=0.9\textwidth]{fig_rr_combined_par_rank.pdf}
  \caption{Rambachan--Roth Sensitivity Analysis: Parental Income Rank}
  \label{fig:rr_par_rank}
  \vspace{1mm}
  \begin{minipage}{0.9\textwidth}
    \footnotesize
    \textit{Notes:} Honest confidence intervals under the smoothness
    ($\Delta_\text{SD}$) and relative magnitude ($\Delta_\text{RM}$)
    restrictions of \textcite{rambachan2023more}. See
    Figure~\ref{fig:rr_par_q1} notes for details. Data:
    \textcite{chetty2020mobility}.
  \end{minipage}
\end{figure}

\subsection{Permutation Inference}
\label{sec:robustness_permutation}

Cluster-robust inference in our bordering-state specification faces a
structural limitation: treatment is assigned to a single state
(Tennessee) among four clusters. With $G_1 = 1$ treated cluster, the
wild cluster bootstrap has no power to reject the null
\parencite{mackinnon2017wild}---the minimum achievable $p$-value under
Rademacher weights is 0.375. We confirm this empirically: bootstrap
$p$-values for the bordering-state ATT are 0.50 (Webb) and 0.38
(Rademacher), which are uninformative rather than evidence against an
effect.

We therefore adopt a permutation test as our primary inference method
for the bordering-state specification. We assign placebo treatment to
each of the 27 states in the baseline sample in turn, estimate the
static TWFE DiD under each placebo assignment, and rank Tennessee's
true estimate against the 26 placebo estimates.

Tennessee's parental income rank effect ($-0.0185$) is the most
extreme of all 27 states (permutation $p = 0.037$, significant at the
5\% level). Tennessee's bottom-quintile share effect ($+0.0154$) is the
second most extreme, exceeded only by Ohio (permutation $p = 0.074$,
marginally significant at the 10\%
level).\footnote{Ohio's large placebo estimate ($+0.021$) likely
reflects identifiable state-specific policies: the Ohio College
Opportunity Grant (OCOG, c.\ 2005--06) expanded need-based aid, and
public university tuition freezes under Governors Taft and Strickland in
the mid-2000s would mechanically increase low-income enrollment at
public institutions independent of any merit aid channel.}
Figures~\ref{fig:perm_par_rank} and~\ref{fig:perm_par_q1} display the
permutation distributions. Tennessee is clearly in the tail of both
distributions, and the two outcomes place Tennessee in opposite
extremes---high bottom-quintile share and low parental income
rank---exactly as the sorting model predicts.

\begin{figure}[H]
  \centering
  \includegraphics[width=0.9\textwidth]{fig_permutation_par_rank.pdf}
  \caption{Permutation Test: Mean Parental Income Rank}
  \label{fig:perm_par_rank}
  \vspace{1mm}
  \begin{minipage}{0.9\textwidth}
    \footnotesize
    \textit{Notes:} Distribution of placebo static TWFE DiD estimates
    obtained by assigning treatment to each of the 27 sample states in
    turn. The vertical line marks Tennessee's true estimate ($-0.0185$),
    which is the most extreme of all 27 states (permutation $p = 0.037$).
    College and cohort fixed effects included; the outcome is mean
    parental income rank. Data: \textcite{chetty2020mobility}.
  \end{minipage}
\end{figure}

\begin{figure}[H]
  \centering
  \includegraphics[width=0.9\textwidth]{fig_permutation_par_q1.pdf}
  \caption{Permutation Test: Bottom-Quintile Share}
  \label{fig:perm_par_q1}
  \vspace{1mm}
  \begin{minipage}{0.9\textwidth}
    \footnotesize
    \textit{Notes:} Distribution of placebo static TWFE DiD estimates
    obtained by assigning treatment to each of the 27 sample states in
    turn. The vertical line marks Tennessee's true estimate ($+0.0154$),
    which is the second most extreme of all 27 states (permutation
    $p = 0.074$). Only Ohio ($+0.021$) produces a larger placebo
    estimate. Data: \textcite{chetty2020mobility}.
  \end{minipage}
\end{figure}

\subsection{Synthetic Control}
\label{sec:robustness_synth}

We implement a synthetic control approach
\parencite{abadie2010synthetic} as a complementary identification
strategy. The synthetic control gap for bottom-quintile share reaches
2.2 percentage points by the 1991 cohort, consistent with the
Callaway--Sant'Anna dynamic path (Figure~\ref{fig:synth_trends_par_q1}).
Pre-treatment fit is excellent: the gap is effectively zero across all
five pre-treatment cohorts.

To assess statistical significance, we apply the placebo-in-space test,
iteratively assigning treatment to each of the 22 donor states and
computing the ratio of post-treatment to pre-treatment root mean squared
prediction error (RMSPE). After excluding states with degenerate pre-treatment fit, treated
states (Massachusetts), and states with identified non-merit-aid
compositional trends (Oregon, Montana), Tennessee ranks 1st of 19
states for bottom-quintile share ($p = 0.053$) and 1st of 20 for
parental income
rank ($p = 0.130$). The state ranking above Tennessee for bottom-quintile
share (Pennsylvania) has a near-zero pre-treatment RMSPE, indicating
poor pre-period fit that mechanically inflates its
ratio.\footnote{Excluding placebo states with pre-treatment RMSPE below
0.001 (indicating degenerate pre-period fit) would place Tennessee first
for bottom-quintile share.}
Figures~\ref{fig:synth_placebo_par_q1}
and~\ref{fig:synth_placebo_par_rank} display the placebo distributions
for both outcomes.

\begin{figure}[H]
  \centering
  \includegraphics[width=0.85\textwidth]{fig_synth_placebo_par_q1.pdf}
  \caption{Synthetic Control Placebo-in-Space: Bottom-Quintile Share}
  \label{fig:synth_placebo_par_q1}
  \vspace{1mm}
  \begin{minipage}{0.85\textwidth}
    \footnotesize
    \textit{Notes:} Each line represents the synthetic control gap
    (actual minus synthetic) for one state, normalized by its
    pre-treatment RMSPE. Tennessee (bold) ranks 1st of 19 states in the
    post/pre RMSPE ratio ($p = 0.053$). Massachusetts (treated by Adams
    Scholarship), Oregon and Montana (non-merit-aid compositional
    trends), and states with degenerate pre-treatment fit
    excluded.\footnote{Oregon's progressive shift is concentrated
    exclusively at community colleges and coincides with sharp tuition
    increases at four-year publics---a within-sector reallocation
    unrelated to merit aid. Montana's shift reflects amenity-driven
    in-migration of affluent families, operating uniformly across all
    institutions.} Data: \textcite{chetty2020mobility}.
  \end{minipage}
\end{figure}

\begin{figure}[H]
  \centering
  \includegraphics[width=0.85\textwidth]{fig_synth_placebo_par_rank.pdf}
  \caption{Synthetic Control Placebo-in-Space: Mean Parental Income Rank}
  \label{fig:synth_placebo_par_rank}
  \vspace{1mm}
  \begin{minipage}{0.85\textwidth}
    \footnotesize
    \textit{Notes:} Each line represents the synthetic control gap
    (actual minus synthetic) for one state, normalized by its
    pre-treatment RMSPE. Tennessee (bold) ranks 1st of 20 states
    ($p = 0.050$). Same exclusions as
    Figure~\ref{fig:synth_placebo_par_q1}. Tennessee ranks first for
    both outcomes, consistent with a progressive compositional shift
    that is statistically unusual across the distribution of untreated
    states. Data: \textcite{chetty2020mobility}.
  \end{minipage}
\end{figure}

Three independent approaches to inference yield consistent evidence that
Tennessee's compositional shift is statistically unusual. Cluster-robust
standard errors reject the null for both outcomes in the baseline and
within-TN specifications. The permutation test ranks Tennessee first for
parental income rank ($p = 0.037$) and second for bottom-quintile share
($p = 0.074$). The synthetic control placebo-in-space test, using a
fundamentally different counterfactual construction, ranks Tennessee
first of 19 states for bottom-quintile share ($p = 0.053$) and first
of 20 for parental income rank ($p = 0.050$). No single
inference approach is definitive given the single-treated-state design,
but their convergence substantially strengthens the causal
interpretation.

\subsection{Summary}

Table~\ref{tab:robustness_summary} consolidates robustness results. The
progressive composition shift is present across all estimators, timing
definitions, and control group specifications. Point estimates for
bottom-quintile share range from 1.27 to 1.64 pp; all are significant at
$p < 0.001$. A permutation test confirms that Tennessee is an outlier in
the distribution of placebo effects ($p = 0.037$ for parental income
rank, $p = 0.074$ for bottom-quintile share). The main vulnerability is
pre-trend sensitivity under the $\Delta_\text{SD}$ restriction, which we
flag as a limitation rather than a resolution.
